% Generated from guochuang-technical-chapters.zh-CN.md; edit the Markdown source.
\chapter{面向大型工程理解的系统框架}
\label{chap:architecture}
\label{chap:challenges}

\section{从功能描述到实现位置的检索问题}

大型工程中的一次功能修改通常始于自然语言需求，最终却需要落实到具体的函数、类型、配置和调用关系。两者之间存在明显的表示差异：需求描述的是业务行为，代码按照语言语法和工程结构组织，而完整行为又可能分散在多个文件甚至不同语言中。例如，“调整上传文件的大小限制”既涉及入口方法，也可能涉及流式读取过程中的累计检查、异常类型和上层错误处理。只返回名称中包含 upload 或 size 的片段，不能充分说明修改应当发生在哪里。

这种差异使代码检索面临两个相互关联的问题。首先，用户往往不知道目标符号的准确名称，单纯的关键词匹配容易遗漏已有实现；其次，即使找到了语义相近的函数，局部正文也未必包含理解和修改该函数所需的依赖。前者要求系统建立从业务职责到代码实体的映射，后者要求系统在命中之后继续恢复实现之间的联系。因此，检索结果除相关源码外，还应列出代码所在的文件和版本、调用的接口，以及暂时无法确定的依赖。开发者据此判断代码是否适用于当前任务，并决定还需查看哪些文件。

工程规模增大后，查找源码和补充依赖所需的时间、内存及模型输入也会增加。若每个任务都重新扫描工程并让模型阅读大量文件，相同结构将被反复解析，模型输入和查询时延也会随工程增长。将全部源码提前切块并向量化虽然能够减少重复扫描，但固定切块会割裂函数和依赖，且高相似度片段可能集中于同一个实现。为减少这部分开销，系统预先解析工程并保存索引；收到查询后，只读取与任务有关的模块、类或函数。这样，每次查询无需重新处理整个仓库。

RECAST 的处理过程分为四步：按源码版本建立索引，根据需求定位模块和符号，补充相关调用与类型定义，最后在模型输入上限内筛选并整理源码。前两步减少重复解析和无关搜索，后两步补足修改所需信息并去除重复内容，供 Agent 制定修改方案和执行验证。

\section{离线建模与在线查询的分层架构}

系统将处理过程划分为离线建模与在线查询两条路径。离线处理保存源码快照，提取文件、符号和依赖记录，再生成模块树、全文索引和向量索引。在线查询读取已发布的索引版本，查找候选代码，补充直接相关的依赖，最后将源码及其位置、版本整理为模型输入，本文称为“上下文”。解析器和数据库负责源码提取与存取，模型主要判断模块职责和候选代码是否符合需求。

\begin{center}\small
\begin{longtable}{P{0.19\textwidth}P{0.33\textwidth}P{0.38\textwidth}}
\toprule
层次 & 核心对象与处理 & 向下游提供的能力 \\
\midrule\endhead
工程接入 & 仓库标识、项目边界、源码快照、内容哈希 & 确定一次分析的输入范围与版本 \\
结构与职责建模 & 声明提取、关系绑定、模块提案、层级校验 & 从文件与符号映射到功能职责 \\
索引与版本存储 & 结构记录、模块结果、全文与向量索引 & 按版本和范围执行局部查询 \\
任务检索 & 多种索引召回、候选融合、工程约束、依赖扩展 & 将自然语言需求定位到源码片段 \\
上下文与开发协同 & 去重、预算控制、未解析与省略说明、Agent 接入 & 为修改、适配和验证提供可追溯输入 \\
\bottomrule
\end{longtable}\end{center}

其中，模块树与依赖图承担不同作用。模块树表达功能范围的包含关系，适合回答“哪个子系统负责这项功能”；依赖图表达调用、引用与绑定关系，适合回答“修改这个入口还涉及哪些实现”。向量索引中的近邻图只负责加速相似度搜索。三者通过模块或符号标识关联。模块归属、代码依赖和向量相似性分别记录，不能用“语义相近”代替“存在调用关系”。

离线与在线分离之后，关键问题转为如何防止分析结果与源码错配。系统为每次读取固定仓库、分析版本和项目范围，并要求符号、模块与源码片段服从同一版本约束。新版本构建期间，已经开始的查询继续读取原来的索引版本；新索引校验并发布后，后续请求再读取新版本。这样，后台更新不会使一次查询前半段读取旧声明、后半段读取新实现。

\section{技术机制之间的协同关系}

代码理解不能只依靠更大的模型输入。结构解析能够稳定提取声明和位置，却难以直接判断业务职责；模型能够解释职责，却可能给出重叠、遗漏或缺乏源码依据的模块划分。本项目因此以结构记录约束模型提案，由模型解释边界、宿主验证归属，再将通过校验的结果写入模块树。模块树的每一次细化都对应可检查的文件集合，而不是仅增加一层自然语言摘要。

有了模块树，检索可以从职责范围缩小到类和函数，但上层判断仍可能有偏差。如果查询只能沿一个模块分支向下搜索，错误路由会使正确实现从候选中消失。因此，模块内检索和全局符号检索共同提供候选：模块摘要提供范围线索，声明和正文检索保留直接命中路径，融合阶段再处理重复及粒度差异。

合并各检索通道的候选结果后，系统还需查找理解这些代码所必需的依赖。扩展不能无限沿图传播，否则公共工具、配置和调用方会迅速占满上下文。系统据此引入任务相关性、关系类型和资源上限，将扩展限制在解释主要实现所需的局部子图；随后合并重叠的源码范围，并按照输入长度限制筛选内容。最终结果包含相关实现、调用关系和源码位置，同时列出未能解析或因长度限制而省略的内容。

\chapter{依赖约束的分层建模与增量索引}
\label{chap:indexing}

\section{多语言源码的统一结构表示}

多语言工程的首要困难，是不同语言前端产生的对象无法直接比较。Java 的类与方法、TypeScript 的导出声明、C/C++ 的函数和注册入口，在语法节点、名称作用域和构建方式上均有差异。如果检索侧直接依赖各语言的原始 AST，就需要为每一种语言重写查询逻辑，也难以把跨语言关系放到同一条路径中。

系统因此将语言前端的解析结果归一化为文件、符号和依赖三类记录。文件保存语言、角色、内容哈希和解析状态；符号保存声明种类、限定名称、签名及源码范围；依赖保存源端、目标端、关系类型、解析状态和对应的源码位置。语法前端负责发现“声明了什么”，后续名称解析或编译器分析负责确定引用所指向的声明，两者共同写入统一模型，查询接口不直接接触解析器内部节点。

符号标识不能仅由行号生成。插入注释或调整排版会改变位置，却未必改变声明本身。现有实现结合文件内 AST 声明标识、限定名称、签名和语言构造符号键，再把行列范围作为版本内的位置属性。源码快照提供实际正文，文件与内容哈希用于核对输入。由此，声明标识、版本位置和实际文本可以分别维护，减少格式变化对关联结果的干扰。

语法解析也不能代替编译语义。重载、条件编译和不同依赖版本可能让同一名称对应多个目标，而缺失 SDK 或构建配置又可能使目标无法确定。系统为关系分别记录已解析、有歧义或未解析状态：只有能够确定源声明和目标声明的关系才建立连接；其余关系记录候选目标或未解析原因，供后续查询和人工核查。目标语言支持因而按“结构可解析、语言内可绑定、跨语言可追踪”逐级建设，而非仅凭文件扩展名宣称完整支持。

\section{大仓解析中的峰值内存与流水线背压}

大型源码建库的资源压力不仅来自输入文件。解析树、符号表、依赖候选、模型输入和待写入记录可能同时驻留内存；当数据库写入速度低于解析速度时，尚未落库的结果会继续累积。即使每个文件都能独立解析，无界的阶段队列仍可能使整体内存失控。因此，优化目标应从“单个解析任务更快”扩展到“各阶段以受控速率协同推进”。

本项目将建库拆为扫描快照、结构解析、关系构建、结果写入与检索索引几个阶段。文件读取与解析采用受限批次，解析工作进程按批次回收，结果再分批写入。当前流水线已经限制了解析资源和写入批量；进一步的背压设计同时约束待处理文件数与字节量，当下游队列达到上限时暂停上游生产，避免大量小文件和少量大文件共享一个不合适的计数阈值。

设阶段 i 的并发数为 p\_i，单个工作单元的平均服务速率为 μ\_i，则稳定流水线的吞吐受最慢阶段限制：

\begin{equation}
X \leq \min_i\left(p_i\mu_i\right).
\end{equation}

增加解析并发只在解析是瓶颈时有效；若写入已饱和，继续增加进程主要扩大队列和内存。峰值内存可进一步分解为：

\begin{equation}
M_{\mathrm{peak}}\approx M_{\mathrm{facts}}+p_{\mathrm{parse}}M_{\mathrm{worker}}+M_{\mathrm{queue}}+M_{\mathrm{write}}.
\end{equation}

其中，四项分别表示驻留结构记录、解析工作进程、阶段队列和写入缓冲的开销。该表达式用于识别资源来源，不能以某个阶段的采样峰值代替整个系统峰值。

解析资源被限制后，结构写入和检索索引生成可能成为新的瓶颈。源码正文、声明信息和向量记录的更新频率与读取方式不同，故应分别保存并按需读取；写入端按批次提交，检索索引生成阶段仅处理需要编码的表示。后续分片与可恢复的任务调度以缩小 M\_facts 和故障重算范围为目标，而不是假定增加线程便能线性提高千万行建库性能。

\section{从目录边界到功能职责的自适应分层}

目录是源码的物理组织方式，功能模块则反映代码共同承担的职责，两者并不总是一致。一个目录可能同时包含多个业务能力，同一能力也可能跨越接口层、服务层和持久化层。直接以目录作为最终模块容易造成职责混杂；反过来，让模型一次读取整个仓库，又会使输入规模和输出校验成本随工程增长。

系统以目录、包、文件角色和已解析依赖生成有限候选组，再由 Agent 判断候选之间的职责关系。每个候选提供文件规模、代表路径、公开接口与源码片段编号，组间依赖汇总成边界线索；模型仅在必要时读取少量代表源码。这样，模型处理的是能够解释的局部范围，宿主仍保有完整的文件集合，可独立核对提案是否遗漏或越界。

对父范围 m 的一次拆分，设文件集合为 F(m)，子范围为 m\_1 至 m\_k。有效拆分要求至少产生两个子范围、子集合不相交，并保持父范围的文件归属完整：

\begin{equation}
k\geq 2,\qquad \bigcup_{i=1}^{k}F(m_i)=F(m),\qquad F(m_i)\cap F(m_j)=\varnothing\ (i\neq j).
\end{equation}

父节点不重复直接拥有子节点的文件，而是聚合其文件、类型和方法统计。无法解析或源码信息不足的文件通过未归属记录说明；校验失败的提案不能直接发布。上述约束把自然语言划分转换为可执行的集合检查，使模型解释与工程完整性相互配合。

划分深度由职责、规模和剩余预算共同决定。若当前范围已经具有单一职责，继续拆分只会增加导航层次和摘要成本，应当停止；若存在可区分的接口与实现边界，则继续细化。输入不足、模型失败或预算耗尽时，节点进入待细化状态，保留现有粗粒度范围及停止原因。小工程可以停在浅层，大工程的不同分支也可以具有不同深度，因而不需要为所有项目强制生成固定六层结构。

自适应划分的调度还受到父子依赖约束。子任务的输入范围只有在父提案通过校验后才能确定，而同一父节点下的独立分支可以并行分析。设计中的调度单位因此是“版本、父范围、候选集合与分析配置”共同标识的任务，解析、模型和向量计算分别使用受限资源池。完成一个分支即可保存其结果，不必等待整棵树结束后统一落盘。

为使并行任务能够恢复，持久记录需要保存输入内容哈希、运行状态、尝试编号和输出哈希。工作进程领取任务后获得本次执行编号，超时重试会生成新的编号；写入结果时同时核对输入版本与执行编号，迟到结果不能覆盖新尝试。父摘要在所需子结果满足条件后聚合，未完成分支继续显式保留。因此，重试可以复用已完成的分支，同时拒绝旧任务迟到的写入，防止同一棵模块树混入不同轮次的分析结果。

\section{构建上下文与跨语言绑定恢复}

统一记录解决了“如何表示关系”，但跨语言关系仍需要找到实际连接点。Java 到 Native 的调用可能经过 JNI 命名约定或动态注册，ArkTS 到 C/C++ 的调用可能经过 N-API 导出，KMP 则通过公共声明与平台实现连接。同名匹配在这些场景下并不充分：导出名、函数签名、注册范围和参与构建的目标必须共同成立，才能形成可信路径。

跨语言分析先识别上层调用入口，再根据导出名、签名和注册项查找底层实现，并记录每一步对应的源码位置。以注册式 Native 调用为例，先从上层声明确定方法名和签名，再查找注册表中对应项，最后解析该项引用的底层函数。图中分别保留声明到注册项、注册项到实现的连接及对应的源码位置。若注册名相同但签名或库范围不同，则保留为不同候选，不能用较高文本相似度消除真实歧义。

构建上下文用于限定候选适用范围。Java 需要考虑依赖类路径和目标版本，C/C++ 需要考虑编译单元、宏与头文件搜索路径，KMP 需要区分公共源集和平台源集。分析结果因此应绑定构建配置哈希；配置变化即使没有修改函数正文，也可能改变关系目标。仅按语法匹配的关系与经过编译器确认的关系分别标记，查询时可以据此判断是否还需核对构建配置。

图查询围绕任务入口提取有限路径：先经过直接调用或绑定边到达实现，再按任务需要查看类型、配置与调用方。跨语言路径可以跨文件和模块，但其每一段都必须保留来源。性能评价同时检查依赖边准确率和关键路径覆盖，因为若一条迁移所需路径中间缺少注册边，即使其余多数边正确，也不足以支撑完整功能追踪。

\section{基于变化类型的增量失效与版本发布}

源码持续更新后，全量重建会重复解析未变化文件，而仅更新改动文件又可能留下失效关联。例如，某个公开方法被删除，调用方正文虽然未变，旧依赖却已不再有效；构建依赖版本变化也可能影响多个符号的绑定。因此，增量更新需要区分“输入是否变化”和“派生结论是否仍成立”。

系统首先比较内容哈希，确定变化文件集合，并复用兼容的解析结果。随后根据变化类型扩展失效范围：方法正文变化主要影响实现片段与所属模块说明，公开签名变化还影响引用该声明的关系，构建配置或注册信息变化则影响依赖其条件的绑定。模块摘要沿包含关系向必要祖先传播，向量缓存依据实际编码文本与模型名称和版本复用。文件路径相同并不足以证明向量可以复用。

设初始变化集合为 A\_0，依赖这些文件或符号生成的记录集合为 Derived(A)，需要更新的模块祖先为 Ancestors(A)，则增量分析反复执行以下扩展，直至不再产生新的失效对象：

\begin{equation}
A_{t+1}=A_t\cup\mathrm{Derived}(A_t)\cup\mathrm{Ancestors}(A_t).
\end{equation}

实现时根据源文件、符号与生成记录之间的反向映射定位受影响的关系，并按正文、接口和构建变化限制传播规则，避免把每次局部修改都扩展成整仓重建。无法安全界定的变化应扩大重算范围，以正确性优先于缓存命中率。

重建结果写入新的分析版本，结构、模块与检索索引分别记录进度。只有必需结果通过完整性和可查询性检查，才切换对外可见版本；失败时保留上一有效版本及失败原因。正在执行的查询继续持有其原有版本，旧结果在不再被读取后回收。增量正确性通过相同最终源码下的全量与增量结果对照验证，再统计重算文件比例、发布时延和资源收益。

\chapter{任务驱动的分层检索与上下文整理}
\label{chap:methods}

\section{检索粒度与模块范围选择}

检索粒度决定返回单个函数、类还是整个模块。函数适合定位局部逻辑，类或接口适合查看类型定义和方法约定，模块与子系统则适合解释职责及较大修改范围。若所有查询只返回固定大小代码块，模块迁移会缺少范围，而局部修复又可能收到过量材料。因此，系统将用户选择的主要结果粒度，与后续补充源码和依赖的粒度分开处理。

请求首先固定仓库、分析版本和工程范围，再选择自动、函数、类、模块或子系统模式。显式粒度具有优先权；当前范围没有相应结果时返回不可用状态，而不是悄悄换成其他类型。自动模式依据可用索引组织多粒度候选，必要的接口、配置与依赖仍可跨粒度补充。粒度选择只改变查询与结果组织，不触发重新解析。

分层检索先利用上层模块的职责摘要定位相关分支，再在这些范围内检索类和函数。其收益来自减少细粒度比较范围，但错误上层判断也可能放大漏检。因此，设计中保留一条不依赖模块路由的全局声明与正文召回通道，与模块范围内的候选共同进入融合。模块树提供方向，全局通道提供直接命中的补充，两者的配额和效果通过任务集调节。

当一项需求涉及多处修改时，查询不能只在单个模块中搜索。例如，“修改上传限制并调整异常返回”同时涉及校验与错误处理，如果排名靠前的候选全是大小校验函数，异常返回逻辑就可能被遗漏。后续将分别识别需求中的校验和错误处理要求，选择各自对应的实现，并保留两者的调用关系。

\section{声明、正文和模块摘要的混合检索}

自然语言和代码标识符具有不同匹配特征。“上传大小限制”可以与语义相近的说明匹配，但具体的 parseRequest、FileUploadBase 或异常名更依赖词法与结构信息。单一向量召回容易忽略精确标识符，单一全文召回又难以跨越不同命名表达。因此，系统对符号声明、源码正文和模块摘要构建多种检索视图，并在每种视图内结合全文与向量通道。

不同通道的原始分值无法直接比较：全文相关度受到词频与长度影响，向量相似度受到模型与编码方式影响。融合阶段使用候选在各通道中的名次形成稳定贡献。对候选 d，其融合得分为：

\begin{equation}
\mathrm{RRF}(d)=\sum_{j:\ d\in L_j}\frac{w_j}{k_0+\mathrm{rank}_j(d)}.
\end{equation}

其中，L\_j 为第 j 个通道的结果列表，名次从 1 开始，w\_j 为通道权重，k\_0 为平滑常数；未命中的通道不贡献分值。当前存储实现使用以 60 为平滑常数的排名融合，任务服务再分别限制声明、片段和摘要候选。

合并后的候选还需对应到具体函数、类或模块。源码片段命中可能只覆盖函数中部，服务根据片段所属的声明映射到函数或类，避免把某段分支误当成完整 API；多种索引命中同一符号时，合并该符号的检索结果并记录各通道的命中情况，而不是重复占据结果位置。模块结果则关联其已发布范围和代表符号，以便从模块检索结果继续读取具体实现。

候选集合形成后再执行较昂贵的判断。工程、语言和目标接口约束先排除明确不适用的结果，然后对有限候选比较任务覆盖、接口兼容性和边界依赖。这样，精排处理的是已经经过结构筛选的集合，而不是整仓源码。当前任务查询流程主要采用排名和规则选择，后续将加入模型对任务适用性的判断，并在相同候选数量限制下比较排序效果，并单独记录模型调用增加的延迟。

\section{在线读取放大与时延控制}

向量搜索本身并非总是在线查询的主要成本。若系统在筛选之前加载大量文档正文，或为了补充一个符号而读取完整结构索引，数据库传输和对象构建就可能超过近邻搜索耗时。随着工程增长，检索时间会更多地消耗在“没有进入最终上下文的数据”上，形成逻辑上的读取放大。

系统先取得数量受限的候选标识，筛选后再读取对应正文。工程范围在查询阶段过滤，符号、片段与摘要通道分别设置上限；候选融合后，再按标识读取符号和源码范围。任务检索依赖局部符号、邻接关系和源码接口，不通过整仓列表拼出一次结果。由此，在线读取量主要受候选与扩展预算约束，而不是直接等于仓库规模。

为观察无效读取，定义一轮查询的读取放大率：

\begin{equation}
A_{\mathrm{read}}=\frac{\mathrm{Bytes}_{\mathrm{read}}}{\mathrm{Bytes}_{\mathrm{delivered}}}.
\end{equation}

其中，分子为查询实际读取的源码字节数，分母为最终返回的源码字节数；没有返回源码的请求单独报告。当前统计不包含数据库磁盘读取量和网络协议开销。该指标不是越接近 1 越好，因为合理的候选筛选需要读取未入选内容，但持续偏高往往提示过滤过晚或重复取数。结合候选命中率与必需代码覆盖率，才能判断减少读取是否牺牲了检索效果。

在串行阶段边界下，查询时延可分解为：

\begin{equation}
T_{\mathrm{query}}=T_{\mathrm{encode}}+T_{\mathrm{recall}}+T_{\mathrm{rank}}+T_{\mathrm{expand}}+T_{\mathrm{compile}}.
\end{equation}

各项分别对应编码、召回、排序、扩展和上下文整理；一个阶段内存在并发时，以该阶段实际墙钟耗时计量，避免重复累加重叠时间。各阶段共享请求截止时间，独立召回可以在受限并发下执行，已计算的查询向量在通道间复用；依赖扩展接近预算时停止，并报告尚未覆盖的范围。耗时改善应来自缩小不必要的工作量和重复计算，而不是在结果中隐藏超时或将未完成请求计为成功。

\section{相关依赖的查找与需求覆盖}

首次检索得到的函数可以作为查找依赖的起点，但其正文未必包含修改所需的全部信息。函数可能使用外部类型、调用内部工具或依赖构建配置，修改接口还可能影响调用方。若沿全部关系展开，公共依赖又会把局部任务扩展到大量无关源码。因此，图扩展需要同时考虑任务、关系方向和资源成本。

系统先读取主要实现，再查询与它直接相连的部分依赖。函数内部修改优先补充被调用实现和类型，公共接口调整需要关注调用方与实现类，跨语言迁移则优先追踪注册绑定和平台适配路径。每次补充内容时记录它与主要实现或任务要求的关系，并保留未解析关系；已访问符号按版本与标识去重，防止环和高连接节点导致重复探索。

为减少只返回同类代码的问题，方案将任务中的不同要求组成集合 U，并为每项要求设置权重 w(u)。对候选代码片段 e，Cov(e) 表示它涉及的需求项；已选片段集合 C 的覆盖得分定义为：

\begin{equation}
F(C)=\sum_{u\in\bigcup_{e\in C}\mathrm{Cov}(e)}w(u).
\end{equation}

需求项可以对应入口实现、目标接口、关键分支、必要依赖或测试。判断某段代码是否涉及一项需求，需要查看其声明、调用关系和具体逻辑；模型生成的说明只能作为分析线索，不能证明代码行为已经符合要求。

这一表示解释了为什么 Top-K 相似度不能替代上下文选择。两个高度相似的函数可能支持同一项需求，加入第二个函数的新增得分很小；一个相似度稍低的异常定义却可能补上尚未覆盖的错误处理要求。选择过程因此需要比较新增覆盖与输入 token 数，优先补充当前结果尚未包含的接口、分支或依赖，而不是只按相似分数继续追加。

\section{输入长度限制下的源码筛选与整理}

上下文选择同时受到 token、文件数和源码行数约束。若只对源码正文估算 token，最终添加标题、版本、关系和省略说明后仍可能超限；若优先堆入所有元数据，又可能使真正需要修改的实现无法进入窗口。因此，长度限制必须作用于最终返回文本，并在来源说明与实际源码之间保留合理空间。

上下文整理模块（代码中的 Context Compiler）先按仓库、版本、路径和内容哈希识别重复片段，再合并相互包含的源码范围。客户端可以提交已收到片段的标识和内容哈希，避免后续请求重复返回。候选加入前尝试渲染完整 Markdown，检查 token、唯一文件数和源码行数；同时为裁剪说明预留空间。最终用 cl100k\_base 对实际导出文本重新计量，保存实际用量及未能包含的内容。

选择源码时，在需求覆盖得分中扣除重复内容的惩罚：

\begin{equation}
\max_{C\subseteq E}\left[F(C)-\lambda\,\mathrm{Redundancy}(C)\right].
\end{equation}

\begin{equation}
\mathrm{Tokens}(C)\leq B,\quad \mathrm{Files}(C)\leq F_{\max},\quad \mathrm{Lines}(C)\leq L_{\max}.
\end{equation}

E 为候选代码片段集合，B 为输入预算，F(C) 为上一节的覆盖得分，Redundancy(C) 衡量重复源码范围，λ 控制冗余惩罚强度。需求覆盖的重复已在集合并集中去除，冗余项进一步约束文本层面的重复。由于每段代码的输入长度和需求覆盖相互影响，直接枚举所有子集需要指数级搜索，在线查询应使用受限候选上的启发式选择。

在线查询采用贪心近似：每轮优先选择单位 token 带来最多新增需求覆盖、且重复较少的片段。实现先以候选独立序列化的增量 token 估计排序成本，接纳时再对完整文本精确计量，避免为每个候选反复编码全部组合。记 J(C) 为覆盖得分减去冗余惩罚，则候选的增量评分为：

\begin{equation}
\mathrm{Score}(e\mid C)=\frac{J(C\cup\{e\})-J(C)}{\max\left(1,\Delta\mathrm{Tokens}(e\mid C)\right)}.
\end{equation}

每轮挑选仍满足长度约束且新增得分较高的代码片段，合并被包含范围，并在新增内容后重新评估受影响候选。任务明确指定的入口或接口优先保留；长度限制不足以容纳这些内容时，明确列出未返回的部分。规模较小的候选集可以用精确枚举作为离线对照，评价启发式与最优值的差距，而不把启发式输出称为全局最优。

整理后的结果包含源码版本、主要实现、相关依赖，以及未能提供的内容及其原因。partial 表示存在未解析、裁剪或其他限制，已返回的内容仍可使用；Agent 根据未解析或省略说明继续查询符号、依赖或源码。补取继续遵守固定版本和已返回片段去重，形成“检索—补齐—修改—验证”的迭代过程。是否提供了足够信息，需要对照人工标注的必需源码和后续修改结果判断；满足输入长度限制并不意味着信息完整。

\chapter{Agent 协同开发与工程实现}
\label{chap:implementation}

\section{HTTP、MCP 与 IDE 的检索接口}

离线索引和在线算法只有进入实际开发流程，才能减少 Agent 重复搜索与阅读代码的成本。不同客户端若各自拼接检索结果，就容易在版本选择、源码范围和省略原因的解释上产生差异。因此，系统使用统一的 ContextPacket 数据结构返回源码、版本、依赖和省略说明。HTTP、MCP 和 IDE 均调用同一检索服务和上下文整理模块。

HTTP 通过 POST /v1/task-search 接收任务请求，返回结构化上下文；MCP 通过 search\_task\_context 返回适合 Agent 阅读的内容，并提供 get\_symbol、get\_dependencies 和 read\_source\_excerpt 等局部工具。一次请求携带明确范围，服务返回对应的版本与来源。Agent 先查看首次检索结果，再根据未解析或省略的内容补查符号、依赖或源码，避免把整个仓库交给模型反复筛选。

浏览器工作台和 VS Code 扩展将这些协议映射为需求输入、工程范围、检索粒度和源码预览。页面保留用户能够理解的任务选择，预算由宿主默认值和服务约束承担。模块树按需分页读取，模块细化状态、解析失败原因和源码声明分别展示；选择模块只改变任务范围和详情，不改变文件中真实的类与方法列表。

系统以不同状态表达建库进度：结构索引可用意味着能够定位声明，模块结果可用意味着已有职责分析，检索索引可用意味着这些结果已进入召回。检索索引生成失败时可以复用已校验提案重试，模型提案失败则保留原因和上一有效结果。状态分离使用户能够判断哪一步需要处理，也为自动恢复提供明确入口。

\section{参考检索、适配规划与多文件生成}

复用迁移进一步要求回答“参考实现如何适合目标工程”。语义相似只能说明功能接近，不能证明接口、类型和外部依赖能够直接复用。例如，参考代码使用同步返回值，而目标接口要求异步执行，逐句翻译即使语法正确，也可能违背目标调用方式。因此，迁移在检索之后增加显式规划，根据参考代码的行为和目标工程的接口要求制定修改步骤。

现有工作区翻译后端接收 Context 与 Spec，由 Analyzer 整理来源到目标的名称映射、依赖处理策略和步骤依赖，再由 Translator 执行多文件修改。每个步骤声明所需文件和前置步骤，目标读写范围由宿主限定。模型只能在已经提供的上下文和允许读取的文件中补充源码和依赖，每个生成步骤保留所参考代码的路径和版本。

多文件修改还需要处理运行中断与用户后续编辑。服务保存修改前内容、预期修改后内容和应用状态，并在写入前持久记录。恢复执行时重新核对文件基线，回退时先检查全部待恢复文件，确认没有被后续修改覆盖，再执行逆向恢复。单个工作区限制并发写入，使计划状态、文件状态与恢复记录保持一致。

编译反馈将计划与生成连接成可执行循环。服务运行宿主配置的编译命令，取得诊断后交给模型修复；完成任务前再次核对通过编译的源码快照是否仍是当前内容。这样能够避免“某一轮编译通过，随后又修改文件却直接结束”的情况。编译结果只覆盖相应语言和工程约束，行为保持还需要下一节的独立验证。

\section{从编译成功到行为验收}

编译器能发现语法和类型问题，却无法证明上传限制、异常映射或数据处理结果符合需求。若用编译成功作为迁移完成标志，错误实现仍可能通过流程。因此，系统分别记录编译检查、行为测试和代码差异审阅结果，完成状态应说明实际通过了哪些检查。

行为验证方案由宿主配置命令和测试范围，模型不能自行替换验收标准。测试用例围绕 Spec 中的正常路径、边界条件和异常行为组织；对迁移任务，可在相同输入下比较参考实现与目标实现的输出及可观察副作用。测试失败时保留失败用例、输出和对应源码版本，将反馈交给后续修复。若只有编译记录，结果继续标记为 compilation-only。

为使行为验证可以审阅，每次运行需要关联任务、计划哈希、目标基线、文件差异和验证命令，并对测试执行前后的源码内容进行核对。只有全部必需检查通过且检查后未再修改对应文件，才进入待审阅状态；用户审阅后再回填目标工程。通过这些记录，可以从最终改动查回使用的参考代码、修改计划和测试结果。

端到端功能测试按固定任务执行上述步骤：登记工程并发布索引，检索参考实现，生成预算内上下文，形成适配计划，执行多文件修改，运行编译与行为测试，最后核对差异和恢复能力。该过程既检查单项技术可调用，也检查它们之间传递的范围、版本和状态是否一致。

\chapter{实验设计与阶段结果}
\label{chap:validation}

\section{验证基线与实验问题}

本章以 chiparon/weichai 的 9-8-v 分支、提交 c3f09d3（2026 年 9 月 8 日 22:35:58，UTC+8）作为现有原型基线。代码用于核对机制接通情况，仓库运行说明及交接材料用于整理既有实验记录。huawei 分支在此基础上实现新增机制，并合入 2026 年 9 月 9 日上游最新的 codex/guochuang-ui 分支（a02e903）。本次以固定样例验证合并后的功能；没有重新执行大仓建库或远程模型实验。以下历史数字均注明所属基线，新实现的验证范围单独说明。

实验围绕四个问题组织：受限建库能否在真实大仓中控制资源；Agent 层级能否形成合法且有意义的职责范围；多粒度检索能否在合理时延内返回真实代码；上下文与开发协同能否提高后续修改的成功率。前三项已有部分工程记录，第四项及完整跨语言能力仍需要独立任务集验证。

\begin{center}\small
\begin{longtable}{P{0.19\textwidth}P{0.33\textwidth}P{0.38\textwidth}}
\toprule
技术环节 & c3f09d3 已有基础 & huawei 接入及后续验证边界 \\
\midrule\endhead
代码建库 & 源码快照、分批解析、版本存储、局部查询 & 并行解析与在途字节背压；资源分片及千万行验证待开展 \\
职责建模 & Agent 提案、自适应层级、归属校验与发布 & 同级分支并行、决策持久化与重试复用；多宿主调度待扩展 \\
语言与关系 & 已接入语言的结构记录、统一关系接口 & 新增目标语言结构前端及保守绑定；完整 SDK 语义待验证 \\
任务检索 & 多种索引的混合检索、粒度约束、有限依赖扩展 & 召回通道并行、查询向量复用、源码读取字节数与时延统计 \\
上下文 & 源码读取、范围去重、最终文本计量 & 按需求词、代码用途及命中符号选择片段；是否包含必需内容待标注评测 \\
开发协同 & 多文件翻译、编译反馈、运行恢复 & 固定行为测试、源码快照核验、HTTP 查询与回滚；UI 与翻译后端仍需联调 \\
\bottomrule
\end{longtable}\end{center}

基线默认支持 Java、TypeScript、JavaScript、Python、C\#、Go 和 Rust。huawei 分支新增 C/C++、Kotlin 与 ArkTS 结构入口，其中 ArkTS 复用 TypeScript 兼容语法，ArkUI 专用语法保留解析诊断。跨语言样例验证 JNI 短导出名、明确 N-API 注册链及同模块 expect/actual 的连接；重载、重复实现或缺少注册的情况保留未解析或歧义。该结果证明有限规则下的路径可追踪，不能替代编译器、SDK 和实际构建配置验证。

\section{百万行建库：解析受限后的写入瓶颈}

既有容量实验采用真实 VS Code src，源码提交为 c5f02db50896ef310d3c909628e651675644fcef，共 7,708 个文件、2,356,695 物理行，包含源码、测试及其他角色，没有通过复制文件扩充规模。结构索引记录 287,471 个符号和 117,420 条语法依赖；解析成功、部分解析和不支持文件分别为 7,691、10 和 7 个。

\begin{center}\small
\begin{longtable}{P{0.19\textwidth}P{0.33\textwidth}P{0.38\textwidth}}
\toprule
阶段或资源 & 既有测量记录 & 统计范围 \\
\midrule\endhead
扫描解析 & 约 385 秒 & 旧扁平模块基线的源码建库 \\
结构写入 & 约 869 秒 & 受本机数据库日志配额影响 \\
检索索引 & 约 784 秒 & 与扫描解析分别计时 \\
解析进程采样峰值 & 约 376 MiB & 父进程与解析子进程合计 \\
全流程主进程峰值 & 约 1.39 GiB & 不含数据库和模型进程 \\
\bottomrule
\end{longtable}\end{center}

这些结果说明，受限解析已能处理百万行真实输入，但建库时间并不主要集中于语法解析。结构写入与检索索引生成均长于扫描解析，继续提高解析并发未必缩短总时间；应优先检查数据库批量写入、索引维护与日志配额影响。两个内存数字的统计范围不同，也不能据此推断整套系统只需要数百 MiB 内存。

已有增量基础支持复用兼容结构结果，并保留全量与增量等价测试。进一步实验将变化分为正文、公开签名、文件增删、构建配置和跨语言注册五类，对照相同最终源码的全量构建结果，检查旧边撤回、端点重绑定和模块摘要更新，再记录重算比例与发布耗时。这样，增量收益才建立在结果一致性的前提上。

\section{Agent 分层：归属覆盖与细化质量的区别}

9 月 8 日的 DeepSeek 大仓记录复用 VS Code 主工程 7,702 个文件的固定索引，生成 72 个层级节点、4 个根范围，其中包含 66 个模块和 6 个子系统。终端深度为 0 至 3，最多 4 层；14 个节点已拆分，1 个为叶节点，57 个仍待细化，文件归属覆盖为 7,702 / 7,702。

模型日志记录 24 次完成的层级请求，其中 16 个有效决策被用于生成模块树，节点中 15 个标记模型来源、57 个保留结构来源。请求数、决策数和节点数分别描述调用、校验与结果，不应混为同一个成功指标。模块建模约 73.842 秒，检索索引生成约 24.139 秒，未重新解析全部源码；完整遍历核对了父子文件、类型和方法聚合一致性。

文件覆盖完整说明每个文件均能在当前范围中找到归属，却不能证明 57 个待细化范围已经具有准确功能边界。小工程进一步说明了两类指标的差异：Java 目标工程 commons-fileupload-java-skeleton 生成 8 个浅层功能模块，覆盖 58 / 61 个文件，一个构建配置不支持解析，两个源码文件解析失败，均保留原因。该结果没有为了覆盖率把源码信息不足的文件强行分配，也没有为了固定深度增加子系统。

后续模块实验采用人工职责标注，分别计算归属完整性、重复归属、职责一致性和边界依赖，并统计模型调用、待细化比例与恢复重算量。并行版本与串行版本使用相同输入及配置比较结果约束和吞吐；中断恢复实验则在父提案发布、子任务运行和检索索引生成阶段注入故障，核对是否复用已完成结果、拒绝迟到写入并保留合法层级。

\section{检索与上下文：源码返回结果与查询时延}

旧大仓基线使用 64 维 hash 向量查询 CodeEditorWidget executeEdits editor text model，约 4.96 秒返回 17 条源码片段，最终文本为 7,813 / 8,000 token，包含 executeEdits 的完整方法。该记录验证了百万行源码中的有界读取与精确计量，但查询包含真实 API 标识符，不能据此推断一般自然语言任务的排序质量。

最新 Agent 模块结果接入检索后的三项记录如下。它们继续使用 hash 向量，任务中包含模块名称和 API，主要检查根据模型生成的模块树能否查到同一版本的实际源码。

\begin{center}\small
\begin{longtable}{P{0.19\textwidth}P{0.33\textwidth}P{0.38\textwidth}}
\toprule
查询粒度 & 记录耗时 & 结果与上下文状态 \\
\midrule\endhead
子系统 & 7.917 秒 & 命中“会话”，6 个文件，partial \\
模块 & 15.537 秒 & 命中“CLI入口”，6 个文件，partial \\
自动 & 8.231 秒 & 返回“会话”等结果，5 个文件，partial \\
\bottomrule
\end{longtable}\end{center}

三项请求均返回实际源码，但模块查询超过 10 秒目标，说明高层结果的代表源码选择与关联读取仍需分析。正式评估应按编码、召回、排序、依赖扩展和上下文整理分阶段记录耗时，同时统计读取数据量，判断延迟来自数据库访问、候选管理还是上下文整理。三个不同请求不构成完整时延分布，不能替代均值、P95 或超时率。

独立中文语义回归使用 Commons FileUpload Java 源码、TypeScript 工程和 384 维 multilingual E5，既有 HTTP 10/10、MCP 6/6 用例通过。这里的通过表示对应源码及协议断言成立，并非大规模任务准确率达到 100\%。语义模型结果与 hash 容量基线分开报告，后续以统一中文任务集标注入口、关键分支和必要依赖，分别计算检索质量、必需代码覆盖率、冗余和补取次数。

\section{代码迁移验证与消融实验}

翻译交接材料记录了真实模型完成 Java 到 TypeScript 的两文件翻译，并在 10 轮调用后编译通过；当时没有行为测试结果，也没有完整页面操作流程证明。该记录支持“多文件生成与编译反馈可以运行”，尚不能支持功能完备度达标。完整实验需在同一任务上继续执行行为用例，并核对测试对应的源码快照与最终写入结果。

对照组采用关键词搜索、单纯向量召回、多种索引的混合检索和同时补充依赖并按需求筛选源码的完整方案，固定模型、任务、候选预算和时间上限。消融实验分别移除模块摘要、依赖扩展、覆盖选择或范围去重，观察缺少哪一环会增加漏检、重复输入或人工修复。检索指标与开发指标共同记录，避免优化相似度名次却没有改善实际修改。

\begin{center}\small
\begin{longtable}{P{0.19\textwidth}P{0.33\textwidth}P{0.38\textwidth}}
\toprule
评价环节 & 主要指标 & 与技术机制的对应关系 \\
\midrule\endhead
建库与更新 & 总耗时、峰值资源、重算比例、发布时延 & 检验流水线和增量失效的收益 \\
模块与关系 & 职责边界、依赖边准确率、路径覆盖、未解析率 & 检验层级及跨语言连接是否可信 \\
检索 & Recall@K、MRR、nDCG、均值与 P95 时延 & 检验融合和路由的质量与成本 \\
上下文 & 必需代码覆盖率、token、重复范围、补取次数 & 检验源码选择是否完整且重复较少 \\
开发结果 & 编译与行为通过率、人工修复量、恢复成功率 & 检验端到端任务是否完成 \\
\bottomrule
\end{longtable}\end{center}

命题要求的检索准确率至少 90\%、平均检索时间 10 秒以内和代码功能完备度超过 90\%，需要在明确口径下分别验收。检索准确率应先定义任务级成功条件和相关性标签，功能完备度应对应需求行为用例；现有少量接口记录与编译记录不足以证明三项指标同时达标。千万行规模实验采用真实工程，单独记录硬件、模型、数据库、并发与冷热缓存条件，使结果能够复现和比较。

\section{新增实现验证与后续实验}

新增机制以固定样例验证输入、输出和失败路径：并行解析受到工作进程数与在途字节的双重限制；同级模块决策乱序返回后仍按稳定顺序组装；重建协调器后可以复用已校验决策；超时模型回复不会写入有效检查点。跨语言注册变化后重新构建关系，并将增量结果与相同最终源码的全量结果比较，检查旧连接是否被清除。

开发流程样例从真实 Java 源文件建库与检索开始，将固定版本的检索结果交给受控工具调用，生成两个目标文件，再运行语法检查与行为断言。另一个样例的初始实现能通过语法检查，却遗漏负数处理；行为测试拒绝完成，修复并重新验证后才记录 behavior-verified。写入后旧结果失效，验收文件禁止修改，差异记录支持重启后回滚。

这些样例使用确定性模型脚本，检验服务协议和完整执行流程，不代表远程模型在未知任务上的成功率。后续仍需完成结构记录与检索索引分片、按变化类型复用模块结果和完整 SDK 语义验证，再以真实大仓及开发任务集分别评价资源成本和修改质量，检查新增机制能否在实际任务中降低开销并提高修改正确率。
